\documentclass[11pt]{article}

\usepackage[margin=1in]{geometry}
\usepackage{amsmath,amssymb}
\usepackage{hyperref}
\usepackage{enumitem}

\hypersetup{
  colorlinks=true,
  linkcolor=blue,
  urlcolor=blue,
  citecolor=blue
}
\emergencystretch=2em

\newcommand{\Tr}{\mathcal T}
\newcommand{\cM}{\mathcal M}
\newcommand{\Z}{\mathbb Z}

\title{A Literature-Implied Negative Answer to the Arbitrary-Spectrum Sampling Problem for \(q>2\)}
\author{Math discovery packet}
\date{July 3, 2026}

\begin{document}
\maketitle

\section{Status}

This is a lightweight \emph{literature-implied answer}, not a new proof.
The source paper is
\[
\text{V. N. Temlyakov, \emph{Sampling discretization of integral norms of the
hyperbolic cross polynomials}, arXiv:2005.05967.}
\]
In Section 4, Open problem 3 asks whether, for every \(q\in(2,\infty)\) and
dimension \(d\), there are constants \(C_i=C_i(q,d)\) and \(c(q,d)\) such
that for every finite \(Q\subset\Z^d\),
\[
  \Tr(Q)\in \cM(m,q,C_2,C_3)
  \quad\text{whenever}\quad
  m\ge C_1 |Q|(\log(2|Q|))^{c(q,d)} .
\]

The supporting paper is
\[
\text{F. Dai and V. N. Temlyakov, \emph{Sampling discretization of integral
norms and its application}, arXiv:2109.09030.}
\]

\section{Identification}

Remark \textup{AR2} in arXiv:2109.09030 gives the obstruction.  Let
\(\Lambda_n=\{k_j\}_{j=1}^n\subset \Z\) be lacunary.  For
\[
  \Tr(\Lambda_n)=\left\{ \sum_{k\in\Lambda_n} c_k e^{ikx} \right\},
\]
the paper recalls the lower bound from KKLT: for every fixed \(p>2\) and
fixed equivalence constants \(C_1,C_2\), the condition
\[
  \Tr(\Lambda_n)\in \cM(m,p,C_1,C_2)
\]
implies
\[
  m\ge c(p,C_1,C_2)n^{p/2}.
\]

Now set \(p=q>2\), \(d=1\), and \(Q=\Lambda_n\).  If Open problem 3 from
arXiv:2005.05967 had a positive answer, then for some constants depending
only on \(q\) one could take
\[
  m \le C n(\log(2n))^{c}
\]
and obtain a valid Marcinkiewicz discretization for all large \(n\).  But
\[
  n(\log n)^c=o(n^{q/2})
  \qquad (q/2>1),
\]
contradicting the lacunary lower bound.  Therefore the arbitrary-spectrum
\(q>2\) statement is false already in dimension one.

\section{Scope}

This answers only Open problem 3 from the source paper.  It does not answer
the near-linear hyperbolic-cross question in Open problem 1, where the
frequency sets have much more structure, and it does not settle Open problem
2 about the logarithmic entropy exponent.

\section{References}

\begin{itemize}[leftmargin=*]
\item V. N. Temlyakov, \emph{Sampling discretization of integral norms of the
hyperbolic cross polynomials}, arXiv:2005.05967.
\item F. Dai and V. N. Temlyakov, \emph{Sampling discretization of integral
norms and its application}, arXiv:2109.09030.
\item The lower bound cited in arXiv:2109.09030 as KKLT, Section D.20.
\end{itemize}

\end{document}
